%%% FIELD proof-main
Fix \(\alpha\in(0,1)\) and a Gaussian tangent geometry \(\mathcal G=(k,g_0,\mathcal A)\).  Write \(P_g\) for the law \(N(g,I_k)\).  By \(\text{def:power-envelope-handle}\), a randomized focal-certificate rule is a Borel function \(\phi:\mathbb R^k\to[0,1]\); its probability of issuing the focal certificate at \(g_0\) is \(E_{g_0}\phi\), and honesty against every opposite-classification tangent \(g\in\mathcal A\) is exactly
\[
 \sup_{g\in\mathcal A}E_g\phi\leq\alpha.
\]
Consequently the largest attainable focal-certification probability is, directly from that definition,
\[
 \beta^{\star}_{\mathcal G}(\alpha)
 =\sup_{\substack{\phi:\mathbb R^k\to[0,1]\ \mathrm{Borel}\\
                    \sup_{g\in\mathcal A}E_g\phi\leq\alpha}}
   E_{g_0}\phi.
\]
This optimization is indexed by the entire set \(\mathcal A\), not merely by its distance from \(g_0\).  We now prove that this distinction is material by calculating two instances of the optimization.

Let \(k=1\), \(g_0=0\), and \(c>0\).  Denote the standard-normal distribution function and density by \(\Phi\) and \(\varphi\), respectively, and let \(p_u(x)=\varphi(x-u)\) be the density of \(P_u=N(u,1)\).

First consider \(\mathcal A_+=[c,\infty)\).  The likelihood ratio at the nearest opposite point is
\[
 r_+(x):=\frac{p_0(x)}{p_c(x)}
 =\exp\{-cx+c^2/2\},
\]
which is strictly decreasing in \(x\).  Put
\[
 a=c+\Phi^{-1}(\alpha),\qquad A_+=(-\infty,a],
\]
so that \(P_c(A_+)=\Phi(a-c)=\alpha\).  Let \(K=r_+(a)>0\).  If \(0\leq\phi\leq1\) is Borel and \(E_c\phi\leq\alpha\), absolute continuity and \(P_c(A_+)=\alpha\) give
\[
 \begin{aligned}
 E_0\phi-P_0(A_+)
 &=\int_{\mathbb R}(\phi-\mathbf 1_{A_+})r_+\,dP_c\\
 &=\int_{\mathbb R}(\phi-\mathbf 1_{A_+})(r_+-K)\,dP_c
   +K\{E_c\phi-P_c(A_+)\}\\
 &\leq0.
 \end{aligned}
\]
Indeed, on \(A_+\), \(\phi-1\leq0\) and \(r_+-K\geq0\); on its complement, \(\phi\geq0\) and \(r_+-K<0\); and the final term is nonpositive by the constraint.  Thus the constraint at \(c\) alone gives the upper bound \(E_0\phi\leq P_0(A_+)\).  Conversely, for every \(u\geq c\),
\[
 P_u(A_+)=\Phi(a-u)\leq\Phi(a-c)=\alpha,
\]
so \(\mathbf 1_{A_+}\) is feasible for all of \(\mathcal A_+\).  Hence
\[
 \beta^{\star}_{\mathcal G_+}(\alpha)
 =P_0(A_+)
 =\Phi\!\left(\Phi^{-1}(\alpha)+c\right).
\]
The displayed signed-integrand inequality also proves uniqueness up to normal-null sets: equality forces \(\phi=1\) where \(x<a\) and \(\phi=0\) where \(x>a\), because \(r_+(x)-K\) is strictly positive on the first region and strictly negative on the second.

Next consider \(\mathcal A_{\pm}=(-\infty,-c]\cup[c,\infty)\).  For any rule \(\phi\) feasible for this geometry, define its reflection symmetrization by
\[
 \psi(x)=\frac{\phi(x)+\phi(-x)}{2}.
\]
Then \(\psi\) is Borel, takes values in \([0,1]\), and symmetry of \(P_0\) gives \(E_0\psi=E_0\phi\).  Moreover, for every \(u\) with \(|u|\geq c\),
\[
 E_u\psi=\frac{E_u\phi+E_{-u}\phi}{2}\leq\alpha,
\]
because both \(u\) and \(-u\) belong to \(\mathcal A_{\pm}\).  Thus symmetrization preserves feasibility and the objective.

Let \(Q=(P_c+P_{-c})/2\).  Its density is
\[
 q(x)=\varphi(x)e^{-c^2/2}\cosh(cx),
 \qquad
 r_{\pm}(x):=\frac{p_0(x)}{q(x)}
 =\frac{e^{c^2/2}}{\cosh(cx)}.
\]
The ratio \(r_{\pm}\) is even and strictly decreasing as a function of \(|x|\).  The function
\[
 F_c(t):=Q([-t,t])=\Phi(t-c)-\Phi(-t-c),\qquad t\geq0,
\]
is continuous, satisfies \(F_c(0)=0\), tends to \(1\) as \(t\to\infty\), and has derivative
\[
 F_c'(t)=\varphi(t-c)+\varphi(t+c)>0.
\]
Therefore there is a unique \(t_{c,\alpha}>0\) such that
\[
 \Phi(t_{c,\alpha}-c)-\Phi(-t_{c,\alpha}-c)=\alpha.
\]
Set \(A_{\pm}=[-t_{c,\alpha},t_{c,\alpha}]\) and \(K_{\pm}=r_{\pm}(t_{c,\alpha})\).  Every symmetrized feasible rule satisfies
\[
 E_Q\psi=\frac{E_c\psi+E_{-c}\psi}{2}\leq\alpha=Q(A_{\pm}).
\]
Repeating the preceding signed-integrand argument with reference measure \(Q\) yields
\[
 \begin{aligned}
 E_0\psi-P_0(A_{\pm})
 &=\int_{\mathbb R}(\psi-\mathbf 1_{A_{\pm}})(r_{\pm}-K_{\pm})\,dQ
   +K_{\pm}\{E_Q\psi-Q(A_{\pm})\}\\
 &\leq0.
 \end{aligned}
\]
Here \(r_{\pm}-K_{\pm}\geq0\) on \(A_{\pm}\) and is negative off that interval, while \(\psi-1\leq0\) on the interval and \(\psi\geq0\) off it.  This proves the upper bound for every feasible \(\phi\), since symmetrization preserves \(E_0\phi\).

It remains to check that the central-interval rule is feasible for the full composite opposite set.  For \(u\geq0\), differentiation gives
\[
 \frac{d}{du}P_u(A_{\pm})
 =-\varphi(t_{c,\alpha}-u)+\varphi(t_{c,\alpha}+u)\leq0,
\]
because \(|t_{c,\alpha}+u|\geq|t_{c,\alpha}-u|\) and the standard-normal density decreases with the absolute value of its argument.  Since \(P_u(A_{\pm})\) is even in \(u\), it is nonincreasing in \(|u|\).  Its value at \(u=c\) is \(Q(A_{\pm})=\alpha\), because the interval is symmetric and hence \(P_c(A_{\pm})=P_{-c}(A_{\pm})=Q(A_{\pm})\).  Therefore \(P_u(A_{\pm})\leq\alpha\) whenever \(|u|\geq c\).  The indicator \(\mathbf 1_{A_{\pm}}\) is thus feasible and attains the upper bound, proving
\[
 \beta^{\star}_{\mathcal G_{\pm}}(\alpha)
 =P_0(A_{\pm})=2\Phi(t_{c,\alpha})-1.
\]

The feasible set for \(\mathcal G_{\pm}\) is contained in the feasible set for \(\mathcal G_+\), so the second value cannot exceed the first.  The inequality is strict.  If equality held, the attaining central-interval rule for \(\mathcal G_{\pm}\) would also attain the one-sided optimum.  By the uniqueness just proved it would then agree, up to normal-null sets, with \(\mathbf 1_{A_+}\).  But the latter is not feasible for \(\mathcal G_{\pm}\), since
\[
 P_{-c}(A_+)
 =\Phi(a+c)
 =\Phi\!\left(\Phi^{-1}(\alpha)+2c\right)
 >\alpha.
\]
This contradiction establishes
\[
 2\Phi(t_{c,\alpha})-1
 <\Phi\!\left(\Phi^{-1}(\alpha)+c\right).
\]

Finally, both geometries have the same nearest-opposite distance,
\[
 \inf_{g\in\mathcal A_+}|g-g_0|
 =\inf_{g\in\mathcal A_{\pm}}|g-g_0|=c,
\]
but their optimal honest focal-certification probabilities are different.  Hence no function of the scalar \(c\) alone can equal \(\beta^{\star}_{\mathcal G}(\alpha)\) for all tangent geometries.  Nor can one rule attain such a putative universal value: any rule attaining the one-sided value must be \(\mathbf 1_{A_+}\) up to null sets, and that rule violates the two-sided honesty constraint.  By \(\text{def:classification-information-distance}\), \(\delta_n(\theta_n)\) records the nearest opposite classification in the information norm; by \(\text{def:power-envelope-handle}\), information whitening turns its finite local limit into the Euclidean nearest-set distance just used while retaining \(\mathcal A\) in \(\mathcal G\).  Thus \(\lim\delta_n(\theta_n)=c\) supplies only the nearest-boundary summary, whereas the explicit optimization and any attaining classifier require the full incident tangent geometry.  This proves every assertion of the theorem.
